\scanpage{145}
\mspage{362}

If we now identify $\widehat{\mathcal G}_{0,3}$ with $\mathrm{Gl}(2,\mathbb Z)^{\prime\wedge}$, and look at the \deleted{two} images of the elements $g,h,f$ and so on in the quotient
\[
 \mathrm{Gl}(2,\widehat{\mathbb Z})^{\pm1\prime}
 =\mathrm{Gl}(2,\widehat{\mathbb Z})^{\pm1}/\pm1
\]
--- which we denote by the same letters, by abuse of language --- \inserted{we} are led to calculate with elements of $\mathrm{Gl}(2,\widehat{\mathbb Z})$, rather than of $\mathrm{Gl}(2,\widehat{\mathbb Z})'$. We \unsure{will choose} $\alpha,\beta$ in $\widehat{\Pi}_{0,3}$ and $\widehat{\widetilde{\Pi}}_{0,3}$ in accordance with (13), hence
\[
 \alpha,\beta\in\mathrm{Gl}(2,\widehat{\mathbb Z})^{\pm1}
 \quad\text{\uline{modulo sign}},
 \tag{33}
\]
but the corresponding elements
\[
 g=\alpha\rho(\alpha)^{-1},\quad
 h=\beta\sigma_\infty(\beta)^{-1}
 \quad\in\mathrm{Sl}(2,\widehat{\mathbb Z})^{\deleted{\unsure{\pm1}}}
 \tag{34}
\]
are determined in $\mathrm{Gl}(2,\widehat{\mathbb Z})^{\deleted{\unsure{\pm1}}}$ \uline{without sign ambiguity}. Likewise
\[
 \overset{\inserted{\text{(case }h=1\text{)}}}{\alpha_0}
 \;\unsure{\text{and}|,}\;
 \overset{\inserted{\text{(case }g=1\text{)}}}{\beta_0}
 \in\mathrm{Gl}(2,\widehat{\mathbb Z})^{\pm1}
 \quad\text{modulo sign},
 \tag{35}
\]
and then
\[
 \begin{aligned}
 \xi&=\alpha_0\rho(\alpha_0)^{-1}
 \inserted{\quad\text{(case }h=1\text{)}},\quad
 \eta=\beta_0\sigma_\infty(\beta_0)^{-1}
 \inserted{\quad\text{(case }g=1\text{)}}
 \in\mathrm{Sl}(2,\widehat{\mathbb Z})
 \end{aligned}
 \tag{36}
\]
are determined without sign ambiguity. We have
\deleted{\unsure{This done|This fixed} \unsure{(34)|(37)}, the elements}
\deleted{the relations}
\[
 \deleted{\unsure{\varphi=} }f\in\mathrm{Sl}(2,\widehat{\mathbb Z})
 \quad\text{modulo \uline{sign}}
 \quad\text{(for }f\in\widehat{\Pi}_{0,3}\text{ fixed)}
 \tag{37}
\]
and we deduce from this
\[
 \varphi=f\sigma_\infty(f)^{-1},\quad
 \psi=f\rho(f)^{-1}\in\mathrm{Sl}(2,\widehat{\mathbb Z})
 \tag{38}
\]
without sign ambiguity.

\deleted{The relations} In the choice of the \deleted{\unsure{representatives}} representatives

\deleted{\unsure{(39)}\ \illegible{1}\ $\unsure{\varphi|\xi}=$}
\margin{144}

\scanpage{146}

of $\varepsilon_0\unsure{\in}\mathrm{Sl}(2,\widehat{\mathbb Z})'$ in $\mathrm{Sl}(2,\widehat{\mathbb Z})$ there is sign ambiguity
\[
 \varepsilon_0=\pm\begin{pmatrix}1&-1\\0&1\end{pmatrix},
\]
\margin{though the sign $-$ would \unsure{hardly} be a \unsure{reasonable} choice --- one likes to see a unipotent element \deleted{\illegible{1}} generating a 1-parameter \unsure{subgroup}\ldots}
but for
\[
 \ell_0=\varepsilon_0^2=\begin{pmatrix}1&-2\\0&1\end{pmatrix}
 \tag{39}
\]
there is no ambiguity of \uline{sign}. This being so, the relations
\[
 \xi=\ell_0^\nu\varphi\qquad\text{(case }h=1\text{)}
 \tag{40}
\]
and
\[
 \eta=\ell_0^{-\nu}\psi\qquad\text{(case }g=1\text{)}
 \tag{41}
\]
are they true without a correcting \uline{sign} \unsure{??|?!}

The case $h=1$ corresponds to the data (in $\widehat{\widetilde{\Pi}}_{0,3}$) of $\xi$ and $\deleted{\unsure{\varphi}}f$, \inserted{and $\nu$,} satisfying
\[
 \xi\rho(\xi)\rho^2(\xi)=1,\qquad
 \xi=\ell_0^\nu
 \bigl(\underbrace{f\sigma_\infty(f)^{-1}}_{\varphi}\bigr),
\]
\deleted{and the equation in $f$ \unsure{is}}
or again to the data
\deleted{$\alpha_0\in\widehat{\widetilde{\Pi}}_{0,3},\ f\in\widehat{\Pi}_{0,3},$}
\[
 \alpha_0\in\widehat{\widetilde{\Pi}}_{0,3},\quad
 f\in\widehat{\Pi}_{0,3},\quad\nu\in\widehat{\mathbb Z}
\]
satisfying the relations
\[
 \underbrace{\alpha_0\rho(\alpha_0)}_{\xi}
 =\ell_0^\nu\underbrace{f\sigma_\infty(f)^{-1}}_{\varphi}
 \tag{42}
\]
in terms of these data we define
\[
 \left\{\begin{aligned}
 \xi&=\alpha_0\rho(\alpha_0)\\
 \varphi&=f\sigma_\infty(f)^{-1}\\
 g&=\sigma_\infty(f)^{-1}\ell_1^\nu\rho(f)
 =(\ell_0^\nu f)^{-1}(\ell_0^\nu\varphi)\rho(\ell_0^\nu f)
 \end{aligned}\right.
 \tag{43}
\]
This recalled, we lift $\alpha_0,f$ to $\mathrm{Gl}(2,\widehat{\mathbb Z})^{\pm1}$,
\margin{145}

\scanpage{147}
\mspage{363}

with sign ambiguity, \deleted{\unsure{but we}} and \deleted{\illegible{2}}\inserted{$\xi$ and $\varphi$} are thus determined in $\unsure{\mathrm{Sl}|\mathrm{Gl}}(2,\widehat{\mathbb Z})$, independently of the choices of signs, \deleted{\illegible{2}} and likewise $g$.

We \unsure{will have|have} therefore in $\mathrm{Sl}(2,\widehat{\mathbb Z})$
\[
 \deleted{\xi=\alpha_0\rho(\alpha_0)=\xi\,\unsure{\varepsilon\ell_0^\nu}}
\]
\[
 \xi=\varepsilon\ell_0^\nu\varphi
 \qquad\text{i.e.}\qquad
 \alpha_0\rho(\alpha_0)=\varepsilon\ell_0^\nu f\sigma_\infty(f)^{-1},
 \tag{44}
\]
\inserted{\uline{in $\mathrm{Sl}(2,\widehat{\mathbb Z})$},}
with a sign $\varepsilon\in\{\pm1\}$ that is well determined, and the question that arises is whether we \unsure{always|always} have $\varepsilon=+1$.

\deleted{\unsure{Defining} $f,\unsure{\rho_0|\rho}$}
\inserted{Let us note that the formula}
\[
 \rho(\ell_0)=\ell_1
\]
is true in $\mathrm{Sl}(2,\widehat{\mathbb Z})$ (without a correcting \uline{sign}) so the \inserted{last} expression for $g$ in (43) is valid in $\mathrm{Sl}(2,\widehat{\mathbb Z})$ (without a \uline{sign}), which shows that $g\in\mathrm{Sl}(2,\widehat{\mathbb Z})$, \deleted{\unsure{is not}} is a $\rho$-transform of $\ell_0^\nu\varphi$, \deleted{\unsure{therefore}} the latter is a $\rho$-transform of $g$.
\deleted{\illegible{1} $g$, In fact (\illegible{3} $g$) the relation}
\inserted{We have}
\[
 g\rho(g)\rho^2(g)=\varepsilon_0.1
 \qquad\varepsilon\in\{\pm1\},
\]
\[
 \deleted{(\ell_0^\nu\varphi)\rho(\ell_0^\nu\varphi).}
\]
and \deleted{likewise} we deduce from this, if $g'=\alpha g\rho(\alpha)^{-1}$
\deleted{\unsure{$\rho^2\in$}} ($\alpha\in\mathrm{Gl}(2,\widehat{\mathbb Z})$), that we have
\[
 g'\rho(g')\rho^2(g')=\varepsilon_0.1
\]
hence in particular for $g'=\ell_0^\nu\varphi\inserted{=\varepsilon\xi\quad\text{(44)}}$ \deleted{hence}, with $\varepsilon\in\{\pm1\}$ \deleted{\illegible{2}} now $\xi$ given by (42) satisfies
\[
 \xi\rho(\xi)\rho^2(\xi)=1
\]
(\uline{N.B.} we do not have $\rho^3=1$ in $\mathrm{Sl}(2,\widehat{\mathbb Z})$, but $\rho^3=-1$, which implies that the inner automorphism \deleted{\unsure{that it}} \unsure{defined|defines} by $\rho$ is of order 3\ldots), hence
\[
 \varepsilon_0=\varepsilon,
\]
\margin{146}

\scanpage{148}

i.e. the question of the validity of the formula
\[
 \xi=\ell_0^\nu\varphi
\]
in $\mathrm{Sl}(2,\widehat{\mathbb Z})$ is equivalent to that of the validity of the formula
\[
 g\rho(g)\rho^2(g)=1,
\]
which we leave open for the moment. Let us note for reference the following formulas \uline{in} \uline{$\unsure{\mathrm{Gl}|\mathrm{Sl}}(2,\widehat{\mathbb Z})$}

\inserted{$\det\alpha_0\in\pm1,\ \det f=1$,}
\[
\begin{gathered}
\text{(Case }h=1\text{)}\\
\left\{
\begin{aligned}
 &\alpha_0,f\in\mathrm{Gl}(2,\widehat{\mathbb Z})^{\pm1}
 \quad \deleted{\unsure{\text{with}}}\ \text{(arbitrary signs, we \uline{choose} them)}\\
 &\left.
 \begin{aligned}
 \xi&\overset{\mathrm{dfn}}=\alpha_0\rho(\alpha_0)^{-1}\\
 \varphi&\overset{\mathrm{dfn}}=f\sigma_\infty(f)^{-1}
 \end{aligned}
 \right|\ \in\mathrm{Sl}(2,\widehat{\mathbb Z})\\
 &\bigl[g\overset{\mathrm{dfn}}=\sigma_\infty(f)^{-1}\ell_1^\nu\rho(f)
 =(\ell_0^\nu f)^{-1}(\ell_0^\nu\varphi)\rho(\ell_0^\nu f)\bigr]\\
 &\text{\unsure{hence|therefore}}\quad
 \xi=\deleted{\unsure{\ell_0^{\nu}}}\cdot\varepsilon\ell_0^\nu\varphi
 \qquad\varepsilon\in\{\pm1\}\\
 &\left[\begin{aligned}
 &g\rho(g)\rho^2(g)=\varepsilon.1\\
 &g=\varepsilon.\alpha\rho(\alpha)^{-1}
 \qquad\alpha=(\ell_0^\nu f)^{-1}\alpha_0=f^{-1}\ell_0^{-\nu}\alpha_0
 \end{aligned}\right]
\end{aligned}
\right.
\end{gathered}
\tag{45}
\]

\margin{It now seems more practical to me to \uline{define} $g$ by $g=\alpha\rho(\alpha)^{-1}$, so as thus to have $g\rho(g)\rho^2(g)=1$, and to \unsure{place} the sign $\varepsilon=\varepsilon_2$ in the formula $g=\varepsilon\sigma_\infty(f^{-1})\ell_1^{\unsure{\nu}}\rho(f)$ --- \unsure{the sign} \illegible{1} \unsure{does not occur} \illegible{3} \unsure{in} $\xi=\varepsilon\ell_0^\nu\varphi$.}
\margin{\uline{N.B.} \unsure{we write} \illegible{1}\quad $\varepsilon=\varepsilon_2$}

Symmetrically
\[
\begin{gathered}
\text{(Case }g=1\text{)}\\
\left\{
\begin{aligned}
 &\beta_0,f\in\mathrm{Sl}(2,\widehat{\mathbb Z})
 \qquad\text{(modulo signs, we \uline{choose} them)}\\
 &\eta\overset{\mathrm{dfn}}=\beta_0\sigma_\infty(\beta_0)^{-1}\\
 &\psi\overset{\mathrm{dfn}}=\unsure{\psi|f}\rho(f)^{-1}\\
 &h\overset{\mathrm{dfn}}=f^{-1}\ell_0^\nu\sigma_\infty(\rho(f))=
\end{aligned}
\right.
\end{gathered}
\tag{46}
\]

\figure{\begin{tikzpicture}[x=1mm,y=1mm,line width=.35pt,line cap=round]
\draw (0,.1) .. controls (.35,-.08) and (.9,-.06) .. (1.3,.04);
\draw (2.2,.22) .. controls (2.45,.34) and (2.65,.25) .. (2.9,.22);
\draw (4.1,.38) -- (5.4,.55);
\draw (6.7,.57) -- (7.45,.87);
\end{tikzpicture}}
\margin{147}

\scanpage{149}
\mspage{364}

I am going to look at the case $h=1$ (cf. \unsure{(45)}), setting
\[
\left\{\begin{aligned}
 \alpha_0&=\begin{pmatrix}a&b\\c&d\end{pmatrix}\\
 f&=\begin{pmatrix}r&s\\t&u\end{pmatrix}
\end{aligned}\right.
\tag{47}
\]
\margin{change of signs in $\alpha_0,\unsure{-f|f}$ is \unsure{inessential}}
\margin{N.B. In the ``arithmetic'' case (i.e. $f$ \unsure{comes|coming} from $\widehat{\Pi}_{0,3}$, $\alpha_0$ from $\widehat{\widetilde{\Pi}}_{0,3}=\widehat{\Pi}_{0,3}.\{1,\tau_\infty\}$, we will have $\alpha_0,f\equiv1\ (2)$ i.e. $a,d,r,u\equiv1\ (2)$, $b,c,s,t\equiv0\ (2)$}
\unsure{if|where} \deleted{\illegible{1}} $a,b,c,d,r,s,t,u\in\widehat{\mathbb Z}$, with
\[
 \boxed{ab-cd=\delta\in\{\pm1\}}
 \tag{48}
\]
\[
 \boxed{ru-st=1}.
 \tag{49}
\]

This \unsure{being set}
\[
 \xi=\alpha_0\rho(\alpha_0)^{-1}
 =\begin{pmatrix}A&B\\C&D\end{pmatrix}
 \tag{50}
\]
\[
 \varphi=f\underbrace{\sigma_\infty(f)^{-1}}_{{}^{t}f}
 =\begin{pmatrix}R&S\\T&\unsure{U|u}\end{pmatrix}
 =\inserted{\unsure{?}}
 \begin{pmatrix}r^2+s^2&rt+su\\\unsure{rt+su}&t^2+u^2\end{pmatrix}
 \tag{51}
\]
so that we will have (cf. \S\,42, formula (78)
\[
\left\{\begin{aligned}
 \delta A&=(a^2+b^2+ab)-(ac+bd+bc)
 &\delta B&=ac+bd+bc\\
 \delta C&=ac+bd+ad-(c^2+d^2+cd)
 &\delta D&=c^2+d^2+cd
\end{aligned}\right.,
\tag{52}
\]
\margin{N.B. we have $\delta=\unsure{\varepsilon_3}$}
and (\S\,42, (90) and 91)
\[
 R=r^2+s^2,\qquad S=T=rt+su,\qquad U=t^2+u^2.
 \tag{53}
\]
\deleted{\illegible{1}}
\deleted{\illegible{2}}\inserted{We will have} the formulas
\[
 C+D\unsure{+|-}B=1
 \tag{54}
\]
\[
 S=T.
 \tag{55}
\]
The essential relation linking $\alpha_0$ and $f$, i.e. linking the ``\unsure{variables}'' $a,b,c,d$ with $u,r,s,t$, is the relation
\[
 \boxed{\xi=\varepsilon\ell_0^\nu\varphi,\qquad\varepsilon\in\{\pm1\}}
 \tag{56}
\]
\margin{``loop relation''}
which is also written $\ell_0^{-\nu}\xi=\varepsilon\varphi$, i.e.
\margin{148}